% ============================================================
%  Capítulo 4 — Pruebas de Hipótesis
% ============================================================
\section{Pruebas de Hipótesis}

\subsection{Formulación del problema}

\begin{defi}[Hipótesis estadística]
  Una \textbf{hipótesis estadística} es una afirmación sobre el parámetro
  $\theta$ (o sobre la distribución de los datos).
  \begin{itemize}
    \item $H_0$: hipótesis nula (afirmación que se desea refutar).
    \item $H_1$ (o $H_a$): hipótesis alternativa.
  \end{itemize}
\end{defi}

\subsection{Errores de tipo I y II}

\begin{center}
\begin{tabular}{lcc}
  \toprule
  & $H_0$ verdadera & $H_0$ falsa \\
  \midrule
  No rechazar $H_0$ & Decisión correcta & Error tipo II ($\beta$) \\
  Rechazar $H_0$    & Error tipo I ($\alpha$) & Decisión correcta \\
  \bottomrule
\end{tabular}
\end{center}

\begin{defi}[Nivel de significación]
  $\alpha = P(\text{rechazar } H_0 \mid H_0 \text{ verdadera})$.
  Valores típicos: $0.01,\, 0.05,\, 0.10$.
\end{defi}

\begin{defi}[Potencia]
  $\pi(\theta) = P(\text{rechazar } H_0 \mid \theta)$.
  Para $\theta \in H_1$ se desea que $\pi(\theta)$ sea grande.
\end{defi}

\subsection{Test de razón de verosimilitudes (TRV)}

\begin{defi}[Estadístico TRV]
  \[
    \Lambda(\mathbf{x}) =
    \frac{\sup_{\theta \in \Theta_0} L(\theta;\mathbf{x})}
         {\sup_{\theta \in \Theta} L(\theta;\mathbf{x})}.
  \]
  Se rechaza $H_0$ cuando $\Lambda(\mathbf{x}) \leq c$ para algún umbral $c$.
\end{defi}

\begin{teo}[Wilks]
  Bajo condiciones de regularidad y si $H_0$ es verdadera,
  \[
    -2\log\Lambda(\mathbf{X}) \dlim \chi^2_r,
  \]
  donde $r = \dim(\Theta) - \dim(\Theta_0)$.
\end{teo}

\subsection{\texorpdfstring{$p$}{p}-valor}

\begin{defi}[$p$-valor]
  El $p$-valor es la probabilidad, bajo $H_0$, de observar un estadístico de
  prueba tan extremo o más que el observado:
  \[
    p = P_{H_0}(T(\mathbf{X}) \geq T(\mathbf{x}_{\text{obs}})).
  \]
  Se rechaza $H_0$ si $p \leq \alpha$.
\end{defi}

\begin{resultado}
  \textbf{Interpretación}: un $p$-valor pequeño es evidencia en contra de
  $H_0$. No mide la probabilidad de que $H_0$ sea verdadera.
\end{resultado}
